% sections/appendix.tex — Appendix
\section{Lifecycle Rules: Quiesce, Terminate, Fail}
\label{app:lifecycle-rules}

For space, the body of \Cref{sec:model:lifecycle} gave the SOS rules
for \textsc{Spawn}, \textsc{Activate}, and \textsc{Process} only.
This appendix gives the three remaining rules, completing the
six-rule lifecycle of \Cref{def:lifecycle-state}.  Each is an
inference rule over configurations \(\langle \Theta, M, t\rangle\)
in the same style as the body rules.

\subsection{Quiesce}
\label{rule:quiesce}

An \Active{} actor transitions to \Draining{} on receipt of a
quiesce signal from the supervisor.  In \Draining{}, the actor's
\textsc{Process} step continues to consume queued messages but the
actor accepts no new messages from the K2K bus; once the inbox is
empty, the supervisor may transition the actor to \Terminated{}
via the \textsc{Terminate} rule below.

\begin{equation}
\inference[\textsc{Quiesce}]
{
  \Theta(a) = \langle \Active{}, \mathit{snap}, \mathit{live},
                       h, r, \tau\rangle
}
{
  \langle \Theta, M, t\rangle
  \trans{\mathsf{quiesce}(a)}
  \langle \Theta', M, t+1\rangle
}
\end{equation}

\noindent where
\(\Theta' = \Theta[a \mapsto
  \langle \Draining{}, \mathit{snap}, \mathit{live},
          h, r, \tau\rangle]\).

The K2K layer enforces the ``no new messages'' contract by
consulting the destination's lifecycle state at the moment of
\textsc{SendOK}: an envelope addressed to a \Draining{} actor is
short-circuited to the dead-letter store with reason
\state{InvalidTarget} because \Cref{rule:k2k-send} includes the
explicit premise \(\Theta(d) = \langle\Active{}, \ldots\rangle\)
and therefore does not fire when the destination is \Draining{}.

\subsection{Terminate}
\label{rule:terminate}

A \Draining{} actor whose inbox is empty transitions to
\Terminated{}, after which the supervisor may release the actor's
slot (back to \Dormant{}).  We give the two-step variant: first
\textsc{Terminate}, then \textsc{Release}.

\begin{equation}
\inference[\textsc{Terminate}]
{
  \Theta(a) = \langle \Draining{}, \mathit{snap}, \mathit{live},
                       h, r, \tau\rangle \\
  Q(a) = \varepsilon
}
{
  \langle \Theta, M, t\rangle
  \trans{\mathsf{terminate}(a)}
  \langle \Theta', M, t+1\rangle
}
\end{equation}

\noindent where
\(\Theta' = \Theta[a \mapsto
  \langle \Terminated{}, \mathit{snap}, \mathit{live},
          h, r, \tau\rangle]\).

\begin{equation}
\inference[\textsc{Release}]
{
  \Theta(a) = \langle \Terminated{}, \ldots\rangle
}
{
  \langle \Theta, M, t\rangle
  \trans{\mathsf{release}(a)}
  \langle \Theta[a \mapsto \bot], M, t+1\rangle
}
\end{equation}

The two-step pattern (terminate, then release) is operationally
useful: between the two steps, the actor's snapshot and final HLC
are still readable for audit purposes.  The supervisor records this
window in its persistent log; the runtime exposes it through
\texttt{GET\ /api/\allowbreak v1/\allowbreak admin/\allowbreak
actor/:id/\allowbreak snapshot} for as long as the operator has not
issued \textsc{Release}.

\subsection{Fail}
\label{rule:fail}

A fault in any non-terminal lifecycle state moves the actor to
\Failed{}.  The supervisor decides next whether to restart
(\Cref{rule:restart}) or to terminate.

\begin{equation}
\inference[\textsc{Fail}]
{
  \Theta(a) = \langle \ell, \mathit{snap}, \mathit{live}, h, r, \tau\rangle \\
  \ell \in \{\Initializing{}, \Active{}, \Draining{}\}
}
{
  \langle \Theta, M, t\rangle
  \trans{\mathsf{fail}(a)}
  \langle \Theta', M, t+1\rangle
}
\end{equation}

\noindent where
\(\Theta' = \Theta[a \mapsto
  \langle \Failed{}, \mathit{snap}, \mathit{live}, h, r, \tau\rangle]\).

The \textsc{Fail} rule does not specify the cause; the runtime
distinguishes timeouts (the actor missed a tick deadline),
explicit panics (the actor's \(\delta\) returned an error), and
device faults (the underlying SM raised an unrecoverable error).
The supervisor's restart policy
(\Cref{sec:supervision:policies}) consults the cause to decide
whether to restart, escalate, or terminate; the formal model
abstracts this into a non-deterministic choice between
\textsc{Restart}~(\Cref{rule:restart}) and a chained
\textsc{Terminate}~/~\textsc{Release} sequence.

\subsection{Composite Properties}
\label{app:composite-properties}

The lifecycle rules---\textsc{Spawn}, \textsc{Activate},
\textsc{Process} from \Cref{sec:model:lifecycle};
\textsc{Quiesce}, \textsc{Terminate}, \textsc{Release},
\textsc{Fail} from \Cref{app:lifecycle-rules}; and
\textsc{Restart} from \Cref{rule:restart}---together produce the
lifecycle state machine (read top-to-bottom, then the three
branches out of \Active{}):

\[
\begin{aligned}
&\Dormant{} \;\xrightarrow{\textsc{Spawn}}\;
  \Initializing{} \;\xrightarrow{\textsc{Activate}}\;
  \Active{} \\[2pt]
&\Active{} \;\xrightarrow{\textsc{Process}}\; \Active{}
  \qquad\text{(self-loop)} \\[2pt]
&\Active{} \;\xrightarrow{\textsc{Quiesce}}\;
  \Draining{} \;\xrightarrow{\textsc{Terminate}}\;
  \Terminated{} \;\xrightarrow{\textsc{Release}}\;
  \Dormant{} \\[2pt]
&\Active{} \;\xrightarrow{\textsc{Fail}}\;
  \Failed{} \;\xrightarrow{\textsc{Restart}}\;
  \Initializing{}
\end{aligned}
\]

The rules are mutually exclusive on the source state: at most one
rule applies to any given configuration for a given actor at a
given step.  This is the basis of the determinism argument that
underwrites the soundness theorem of \Cref{thm:lifecycle-soundness}:
since each rule has a unique target state per source state, the
configuration's evolution is determined (modulo the supervisor's
non-deterministic policy choices) by the trace of supervisor
actions and message arrivals.
